% =====================================================================
%  ARTIGO — Álgebra Linear e Geometria Analítica
%  Compilar: pdflatex artigo.tex (duas vezes)  — ou enviar ao Overleaf
%  >>> PREENCHA OS NOMES DOS INTEGRANTES E O LINK DO REPOSITÓRIO ABAIXO <<<
% =====================================================================
\newcommand{\IntegranteUm}{Nome Completo do Integrante 1}
\newcommand{\IntegranteDois}{Nome Completo do Integrante 2}
\newcommand{\IntegranteTres}{Nome Completo do Integrante 3}
\newcommand{\Instituicao}{UniCatólica RN -- Curso de Engenharia de Software}
\newcommand{\Repositorio}{https://github.com/USUARIO/algebra-linear-desempenho-academico}

\documentclass[12pt,a4paper]{article}
\usepackage[T1]{fontenc}
\usepackage[utf8]{inputenc}
\usepackage{mathptmx}                 % Times New Roman (texto e matemática)
\usepackage[a4paper,top=3cm,left=3cm,bottom=2cm,right=2cm]{geometry}
\usepackage{setspace}
\usepackage{indentfirst}
\usepackage{amsmath,amssymb}
\usepackage{graphicx}
\usepackage{booktabs}
\usepackage{float}
\usepackage[font=small,labelsep=endash,justification=centering]{caption}
\usepackage{titlesec}
\usepackage{enumitem}
\usepackage{xcolor}
\usepackage[hidelinks]{hyperref}
\usepackage{url}
\usepackage{fancyhdr}
\setlength{\headheight}{14.5pt}
\renewcommand{\topfraction}{0.85}\renewcommand{\bottomfraction}{0.6}\renewcommand{\textfraction}{0.12}\renewcommand{\floatpagefraction}{0.8}
\renewcommand{\ttdefault}{pcr}

% ---------- ABNT: rótulos em português, recuo, espaçamento ----------
\renewcommand{\figurename}{Figura}
\renewcommand{\tablename}{Tabela}
\setlength{\parindent}{1.25cm}
\onehalfspacing
\hyphenpenalty=10000 \exhyphenpenalty=10000 \emergencystretch=3em
\setlength{\abovedisplayskip}{6pt}\setlength{\belowdisplayskip}{6pt}
% Seções conforme NBR 6024: primária em CAIXA ALTA negrito; secundária negrito
\titleformat{\section}{\normalfont\bfseries\MakeUppercase}{\thesection}{0.5em}{}
\titleformat{\subsection}{\normalfont\bfseries}{\thesubsection}{0.5em}{}
\titlespacing*{\section}{0pt}{12pt}{6pt}
\titlespacing*{\subsection}{0pt}{10pt}{4pt}
% Paginação no canto superior direito (NBR 14724)
\pagestyle{fancy}\fancyhf{}\fancyhead[R]{\small\thepage}
\renewcommand{\headrulewidth}{0pt}
% Fonte da figura/tabela (abaixo, tamanho menor)
\newcommand{\fonte}[1]{\par\vspace{2pt}{\footnotesize Fonte: #1\par}}
\newcommand{\virg}{{,}}
\newcommand{\R}{\mathbb{R}}

\begin{document}
\thispagestyle{fancy}
% ============================ ELEMENTOS PRÉ-TEXTUAIS ============================
\begin{center}
\begin{singlespace}
{\bfseries\large ANÁLISE DO DESEMPENHO ACADÊMICO POR MEIO DE OPERAÇÕES MATRICIAIS: UMA APLICAÇÃO DA ÁLGEBRA LINEAR\par}
\vspace{4pt}
{\itshape Analysis of academic performance through matrix operations: an application of linear algebra\par}
\vspace{10pt}
\IntegranteUm\footnote{Discente. \Instituicao.}\\
\IntegranteDois\footnotemark[1]\\
\IntegranteTres\footnotemark[1]\par
\end{singlespace}
\end{center}

\begin{singlespace}
\noindent\textbf{RESUMO}\par\vspace{3pt}
\noindent Este artigo investiga como operações matriciais podem auxiliar a análise do desempenho acadêmico, permitindo identificar disciplinas de maior dificuldade, comparar rendimentos e calcular indicadores de forma sistemática. Foram utilizadas as notas reais de 382 estudantes de duas escolas públicas portuguesas que cursaram Matemática e Português, disponíveis na base \textit{Student Performance} do repositório UCI. As notas foram organizadas na matriz $N \in \R^{382\times 6}$, cujas linhas representam estudantes e cujas colunas representam os três períodos letivos de cada disciplina. Mudança de escala por multiplicação escalar, soma de matrizes, produto matriz-vetor com o vetor de pesos $w=(0{,}3;\,0{,}3;\,0{,}4)^T$, produto por uma matriz bloco-diagonal, transposição, matriz de covariância e médias por grupo obtidas por $(E^TE)^{-1}E^TF$ foram implementados manualmente em Python e conferidos com a biblioteca NumPy. Os resultados indicam que Matemática é a disciplina de maior dificuldade: média ponderada de 5,31 contra 6,16 em Português, desvio-padrão de 1,92 contra 1,28 e taxa de aprovação de 57,1\% contra 85,9\%, além de queda média de 0,24 ponto entre o primeiro e o terceiro período. A correlação entre as disciplinas é moderada ($r=0{,}58$), e o histórico de reprovações é o fator que mais diferencia os grupos. A análise de sensibilidade mostrou que a escolha dos pesos altera as médias em menos de 0,05 ponto, o que confere robustez ao diagnóstico.

\vspace{4pt}
\noindent\textbf{Palavras-chave:} álgebra linear; matrizes; desempenho acadêmico; multiplicação matricial; análise de dados educacionais.

\vspace{8pt}
\noindent\textbf{ABSTRACT}\par\vspace{3pt}
\noindent This paper uses matrix operations to analyze the real grades of 382 Portuguese secondary students in Mathematics and Portuguese. Grades form a $382\times 6$ matrix; scalar multiplication, matrix addition, weighted matrix-vector products, block-diagonal products, covariance matrices and group means computed as $(E^TE)^{-1}E^TF$ were implemented from their definitions in Python. Mathematics is the most difficult subject (mean 5.31 vs. 6.16; pass rate 57.1\% vs. 85.9\%), and results are robust to the choice of weights.

\vspace{4pt}
\noindent\textbf{Keywords:} linear algebra; matrices; academic performance; matrix multiplication.
\end{singlespace}

% ================================ INTRODUÇÃO ================================
\section{Introdução}

Escolas e universidades produzem, a cada período letivo, grandes volumes de notas. Quando esses dados são examinados apenas aluno por aluno, perguntas importantes ficam sem resposta: qual disciplina concentra mais dificuldade? O rendimento melhora ou piora ao longo do ano? Estudantes que vão mal em uma disciplina tendem a ir mal em outra? A área de mineração de dados educacionais mostra que respostas sistemáticas a essas perguntas apoiam decisões pedagógicas, como reforço escolar e acompanhamento de estudantes em risco (ROMERO; VENTURA, 2010, p.~601).

Uma tabela de notas é, por natureza, uma matriz: cada linha corresponde a um estudante e cada coluna a uma avaliação. Essa observação simples permite que os indicadores usuais (médias, notas ponderadas, taxas de aprovação, comparações entre grupos e correlações) sejam expressos como operações da Álgebra Linear, como transposição, soma, multiplicação por escalar e produto matricial (BOLDRINI \textit{et al.}, 1980; ANTON; RORRES, 2012). A vantagem dessa formulação é que uma única expressão, como $F=NW$, calcula simultaneamente o indicador de todos os estudantes e de todas as disciplinas, com regras de compatibilidade de dimensões que ajudam a evitar erros.

O problema de pesquisa deste trabalho é: \textit{como a utilização de operações matriciais pode auxiliar na análise do desempenho acadêmico de estudantes, permitindo identificar disciplinas com maior dificuldade, comparar rendimentos e calcular indicadores de desempenho de maneira sistemática?}

O objetivo geral é desenvolver uma aplicação em Python, apoiada em conceitos de Álgebra Linear, para analisar dados acadêmicos reais por meio de matrizes, vetores e operações matriciais. Os objetivos específicos são: representar as notas em uma matriz estudantes $\times$ avaliações; calcular médias por estudante e por disciplina; aplicar a multiplicação matricial com vetores de pesos; identificar a disciplina de menor rendimento; demonstrar passo a passo as operações utilizadas; e analisar criticamente os resultados. Utilizou-se a base pública \textit{Student Performance} (CORTEZ; SILVA, 2008; CORTEZ, 2014), com notas de estudantes do ensino secundário de Portugal. O código, os dados e os resultados estão disponíveis no repositório do projeto\footnote{Disponível em: \url{\Repositorio}.}.

% ========================== FUNDAMENTAÇÃO TEÓRICA ==========================
\section{Fundamentação teórica}

\subsection{Matrizes e vetores}

Uma matriz $A$ de ordem $m\times n$ é uma tabela de números reais com $m$ linhas e $n$ colunas, cujo elemento da linha $i$ e coluna $j$ é denotado por $a_{ij}$ (BOLDRINI \textit{et al.}, 1980; STEINBRUCH; WINTERLE, 1987). Um vetor coluna de $\R^n$ é uma matriz $n\times 1$. Neste trabalho, a linha $i$ de uma matriz de notas descreve o estudante $i$, e a coluna $j$ descreve uma avaliação; assim, o vetor $\mathbf{1}=(1,1,\dots,1)^T$ e os vetores de pesos são os instrumentos que \emph{resumem} linhas ou colunas.

\subsection{Operações matriciais}

Dadas $A,B\in\R^{m\times n}$ e um escalar $k\in\R$, a soma e a multiplicação por escalar são definidas elemento a elemento: $(A+B)_{ij}=a_{ij}+b_{ij}$ e $(kA)_{ij}=k\,a_{ij}$. A transposta $A^T\in\R^{n\times m}$ é definida por $(A^T)_{ij}=a_{ji}$, ou seja, troca linhas por colunas (ANTON; RORRES, 2012).

O produto de $A\in\R^{m\times n}$ por $B\in\R^{n\times p}$ só existe quando o número de colunas de $A$ é igual ao número de linhas de $B$, e resulta em $AB\in\R^{m\times p}$ com
\begin{equation}
(AB)_{ik}=\sum_{j=1}^{n} a_{ij}\,b_{jk}.
\label{eq:produto}
\end{equation}
Em particular, se $v\in\R^{n}$, o produto $Av\in\R^m$ tem componentes $(Av)_i=\sum_j a_{ij}v_j$. Lay (2013) interpreta esse produto como uma \emph{combinação linear das colunas de $A$} com coeficientes dados por $v$; quando $v$ tem componentes não negativas que somam 1, cada $(Av)_i$ é uma \emph{média ponderada} da linha $i$. Essa é a interpretação central deste trabalho.

\subsection{Indicadores estatísticos em forma matricial}

Médias, variâncias e correlações também podem ser escritas com matrizes (JOHNSON; WICHERN, 2007). Para $A\in\R^{m\times n}$, o vetor de médias das colunas e o vetor de médias das linhas são
\begin{equation}
\mu=\frac{1}{m}A^T\mathbf{1}_m, \qquad r=\frac{1}{n}A\,\mathbf{1}_n .
\label{eq:medias}
\end{equation}
Subtraindo de cada coluna a sua média obtém-se a matriz centralizada $X_c=A-\mathbf{1}_m\mu^T$, e a matriz de covariância é
\begin{equation}
C=\frac{1}{m-1}X_c^T X_c \in \R^{n\times n}.
\label{eq:cov}
\end{equation}
$C$ é simétrica; sua diagonal contém as variâncias e os demais elementos, as covariâncias. A correlação entre as colunas $j$ e $k$ é $r_{jk}=c_{jk}/\sqrt{c_{jj}c_{kk}}$, que em forma matricial se escreve $R=D^{-1/2}CD^{-1/2}$, com $D=\operatorname{diag}(C)$.

Por fim, para comparar grupos (por exemplo, escolas), usa-se uma matriz indicadora $E\in\R^{m\times g}$, com $e_{ik}=1$ se o estudante $i$ pertence ao grupo $k$ e $0$ caso contrário. Como cada estudante pertence a um único grupo, $E^TE$ é diagonal e contém o tamanho de cada grupo, e
\begin{equation}
G=(E^TE)^{-1}E^TA
\label{eq:grupos}
\end{equation}
fornece, na linha $k$, a média das colunas de $A$ no grupo $k$. Essa expressão é a mesma solução de mínimos quadrados do modelo linear $A\approx EG$ (STRANG, 2016), o que mostra que comparar médias de grupos é um problema de Álgebra Linear.

% =============================== METODOLOGIA ===============================
\section{Metodologia}

\subsection{Dados}

Trata-se de uma pesquisa aplicada, quantitativa, com dados secundários. A base \textit{Student Performance} (CORTEZ, 2014) foi coletada por Cortez e Silva (2008) em duas escolas públicas portuguesas (Gabriel Pereira e Mousinho da Silveira) e contém 395 estudantes de Matemática e 649 de Português. Para cada disciplina há três notas, de 0 a 20: G1 e G2 (1º e 2º períodos) e G3 (3º período, que é a nota final do ano). Os estudantes que cursaram as duas disciplinas foram identificados pela coincidência de 13 atributos, conforme o procedimento indicado pelos autores da base, o que resultou em $m=382$ estudantes, numerados de E001 a E382.

\subsection{Modelagem matricial}

As notas foram organizadas na matriz $A\in\R^{382\times 6}$, cujas colunas são, nesta ordem, Matemática P1, P2, P3 e Português P1, P2, P3. Para adotar a escala brasileira de 0 a 10, aplicou-se a multiplicação por escalar $N=\tfrac{1}{2}A$. As cinco primeiras linhas de $N$ são:
\begin{equation}
N=\frac{1}{2}A=
\begin{bmatrix}
2{,}5 & 3{,}0 & 3{,}0 & 0{,}0 & 5{,}5 & 5{,}5\\
2{,}5 & 2{,}5 & 3{,}0 & 4{,}5 & 5{,}5 & 5{,}5\\
3{,}5 & 4{,}0 & 5{,}0 & 6{,}0 & 6{,}5 & 6{,}0\\
7{,}5 & 7{,}0 & 7{,}5 & 7{,}0 & 7{,}0 & 7{,}0\\
3{,}0 & 5{,}0 & 5{,}0 & 5{,}5 & 6{,}5 & 6{,}5\\
\vdots & \vdots & \vdots & \vdots & \vdots & \vdots
\end{bmatrix}
\begin{matrix}\leftarrow \text{E001}\\ \leftarrow \text{E002}\\ \leftarrow \text{E003}\\ \leftarrow \text{E004}\\ \leftarrow \text{E005}\\ \phantom{\vdots}\end{matrix}
\label{eq:N}
\end{equation}

A matriz é $382\times 6$: o elemento $n_{ij}$ é a nota do estudante $i$ na avaliação $j$. Separando as colunas de cada período obtêm-se as matrizes $P_1,P_2,P_3\in\R^{382\times 2}$ (colunas: Matemática e Português). A média aritmética dos períodos é dada por soma de matrizes e multiplicação por escalar, $\bar{M}=\tfrac{1}{3}(P_1+P_2+P_3)$, e a evolução ao longo do ano, pela diferença $\Delta=P_3-P_1$.

\textbf{Vetor de pesos.} Adotou-se, como parâmetro do estudo, o vetor $w=(0{,}3;\ 0{,}3;\ 0{,}4)^T$, que atribui peso maior ao 3º período por ele corresponder à avaliação final do ano. Os pesos não são oficiais das escolas: são não negativos, somam 1 (o que mantém o resultado na escala 0--10) e podem ser alterados por linha de comando. Sendo $N_{\text{Mat}}\in\R^{382\times 3}$ as três colunas de Matemática, a nota ponderada é $n_{\text{Mat}}=N_{\text{Mat}}\,w$. As dimensões são compatíveis porque $N_{\text{Mat}}$ tem 3 colunas e $w$ tem 3 componentes; o resultado é um vetor de $\R^{382}$, uma nota por estudante. Para calcular as duas disciplinas de uma só vez, usou-se a matriz bloco-diagonal $W\in\R^{6\times 2}$:
\begin{equation}
F=NW,\qquad
W=\begin{bmatrix} w & \mathbf{0}\\ \mathbf{0} & w\end{bmatrix}=
\begin{bmatrix}
0{,}3 & 0\\ 0{,}3 & 0\\ 0{,}4 & 0\\ 0 & 0{,}3\\ 0 & 0{,}3\\ 0 & 0{,}4
\end{bmatrix},
\qquad F\in\R^{382\times 2}.
\label{eq:F}
\end{equation}
Pela Equação~(\ref{eq:produto}), $F=0{,}3P_1+0{,}3P_2+0{,}4P_3$, ou seja, $F$ é uma combinação linear das matrizes de período. O índice global de cada estudante é $g=Fp$, com $p=(0{,}5;\ 0{,}5)^T$ (pesos iguais para as disciplinas).

\textbf{Indicadores.} As médias por disciplina foram obtidas por $\mu_F=\frac{1}{m}F^T\mathbf{1}$ (Equação~\ref{eq:medias}); a dispersão e a associação entre disciplinas, pela matriz de covariância (Equação~\ref{eq:cov}) de $F$ e pela matriz de correlação $6\times 6$ de $N$. A aprovação foi codificada na matriz indicadora $I\in\R^{382\times 2}$, com $I_{ij}=1$ se $f_{ij}\ge 5{,}0$ (equivalente a 10 em 20, mínimo de aprovação no sistema português), de modo que a taxa de aprovação é $\frac{1}{m}I^T\mathbf{1}$. As comparações entre escolas, tempo de estudo semanal e número de reprovações anteriores foram feitas com a Equação~(\ref{eq:grupos}).

\textbf{Análises de robustez.} Foram calculados três cenários de pesos ao mesmo tempo, empilhando os vetores nas colunas de $W_c\in\R^{3\times 3}$ e calculando $N_{\text{Mat}}W_c\in\R^{382\times 3}$. Também se repetiram os cálculos excluindo os estudantes com nota zero no 3º período.

\subsection{Implementação}

O projeto foi escrito em Python 3. As operações (transposta, soma, escalar, produto matriz-vetor, produto matricial, covariância, correlação e médias por grupo) foram implementadas \emph{manualmente} no módulo \texttt{algebra\_matricial.py}, diretamente a partir das definições da Seção~2; por exemplo, o produto matricial é um somatório duplo que reproduz a Equação~(\ref{eq:produto}). A biblioteca pandas (McKINNEY, 2010) foi usada apenas para ler os arquivos CSV, a matplotlib (HUNTER, 2007) para os gráficos, e a NumPy (HARRIS \textit{et al.}, 2020) somente nos testes automatizados, que conferem cada função manual contra o resultado da biblioteca. Os nove testes passaram, e a diferença máxima entre as três formas de calcular $F$ (produtos separados, $NW$ e combinação linear dos $P_k$) foi de $1{,}8\times 10^{-15}$, ou seja, erro de arredondamento. O programa completo é executado com \texttt{python codigo/main.py} em cerca de 1 segundo.

% ========================= RESULTADOS E DISCUSSÃO =========================
\section{Resultados e discussão}

\subsection{Demonstração passo a passo: produto matriz-vetor}

Para ilustrar a operação central, toma-se a submatriz formada pelos quatro primeiros estudantes e pelas notas de Matemática (colunas 1 a 3 de $N$), multiplicada pelo vetor de pesos $w$:

\textbf{1. Matriz e vetor originais.}
\[
N_{\text{Mat}}^{(4)}=\begin{bmatrix}
2{,}5 & 3{,}0 & 3{,}0\\ 2{,}5 & 2{,}5 & 3{,}0\\ 3{,}5 & 4{,}0 & 5{,}0\\ 7{,}5 & 7{,}0 & 7{,}5
\end{bmatrix}_{4\times 3},
\qquad
w=\begin{bmatrix}0{,}3\\0{,}3\\0{,}4\end{bmatrix}_{3\times 1}.
\]

\textbf{2. Operação.} Produto matriz-vetor $n=N_{\text{Mat}}^{(4)}w$. A matriz tem 3 colunas e $w$ tem 3 linhas, logo o produto existe e tem ordem $4\times 1$.

\textbf{3. Cálculo de cada elemento} ($n_i=a_{i1}w_1+a_{i2}w_2+a_{i3}w_3$):
\begin{align*}
n_1 &= 2{,}5(0{,}3)+3{,}0(0{,}3)+3{,}0(0{,}4)=0{,}75+0{,}90+1{,}20=2{,}85\\
n_2 &= 2{,}5(0{,}3)+2{,}5(0{,}3)+3{,}0(0{,}4)=0{,}75+0{,}75+1{,}20=2{,}70\\
n_3 &= 3{,}5(0{,}3)+4{,}0(0{,}3)+5{,}0(0{,}4)=1{,}05+1{,}20+2{,}00=4{,}25\\
n_4 &= 7{,}5(0{,}3)+7{,}0(0{,}3)+7{,}5(0{,}4)=2{,}25+2{,}10+3{,}00=7{,}35
\end{align*}

\textbf{4. Resultado.} $n=(2{,}85;\ 2{,}70;\ 4{,}25;\ 7{,}35)^T$. Os estudantes E001, E002 e E003 ficaram abaixo de 5,0 em Matemática e E004 foi aprovado. Repetindo a conta para Português obtêm-se $3{,}85$; $5{,}20$; $6{,}15$ e $7{,}00$, e o índice global $g=Fp$ desses estudantes é $3{,}35$; $3{,}95$; $5{,}20$ e $7{,}175$. Os valores coincidem com os do arquivo \texttt{notas\_ponderadas\_e\_indice.csv} gerado pelo programa. O caso de E001 merece atenção: a nota zero em Português no 1º período, seguida de 5,5 nos dois períodos seguintes, sugere ausência na avaliação, e não desconhecimento; a ponderação, porém, a trata como nota real e reduz o índice do estudante.

\subsection{Disciplina de maior dificuldade}

A Tabela~\ref{tab:indicadores} reúne os indicadores calculados com $F=NW$.

\begin{table}[htbp]
\centering
\caption{Indicadores por disciplina (notas ponderadas, escala 0--10, $m=382$)}
\label{tab:indicadores}
\begin{singlespace}\small
\begin{tabular}{lcc}
\toprule
Indicador & Matemática & Português\\
\midrule
Média ponderada $\mu_F=\frac{1}{m}F^T\mathbf{1}$ & 5,31 & 6,16\\
Desvio-padrão $\sqrt{c_{jj}}$ & 1,92 & 1,28\\
Taxa de aprovação $\frac{1}{m}I^T\mathbf{1}$ & 57,1\% & 85,9\%\\
Variação média P3 $-$ P1 (matriz $\Delta$) & $-0{,}24$ & $+0{,}20$\\
Estudantes com nota zero no P3 & 39 & 5\\
\bottomrule
\end{tabular}
\end{singlespace}
\fonte{elaborada pelos autores a partir dos resultados do código (2026).}
\end{table}

Matemática é a disciplina de maior dificuldade por todos os critérios: média 0,85 ponto menor, taxa de reprovação de 42,9\% (contra 14,1\% em Português) e maior dispersão, o que indica uma turma mais desigual. A Figura~\ref{fig:periodos} mostra ainda trajetórias opostas: a média de Matemática cai de 5,43 para 5,19 ao longo do ano, enquanto a de Português sobe de 6,06 para 6,26. Do total, 203 estudantes foram aprovados nas duas disciplinas e 39 foram reprovados em ambas.

\begin{figure}[htbp]
\centering
\begin{minipage}[t]{0.48\textwidth}\centering
\caption{Média da turma por período e disciplina}
\label{fig:periodos}
\includegraphics[width=\textwidth]{graficos/fig1_media_por_periodo.png}
\fonte{elaborada pelos autores (2026).}
\end{minipage}\hfill
\begin{minipage}[t]{0.48\textwidth}\centering
\caption{Distribuição das notas ponderadas ($F=NW$)}
\label{fig:dist}
\includegraphics[width=\textwidth]{graficos/fig2_distribuicao_notas_ponderadas.png}
\fonte{elaborada pelos autores (2026).}
\end{minipage}
\end{figure}

A Figura~\ref{fig:dist} mostra a distribuição das notas ponderadas. Em Matemática há uma cauda de notas muito baixas (entre 0,5 e 3), praticamente ausente em Português.


\subsection{Relação entre disciplinas e entre períodos}

A matriz de covariância das notas ponderadas, calculada por $C_F=\frac{1}{m-1}X_c^TX_c$, foi
\[
C_F=\begin{bmatrix} 3{,}6903 & 1{,}4203\\ 1{,}4203 & 1{,}6427\end{bmatrix}.
\]
Sua diagonal fornece as variâncias (e, portanto, os desvios $\sqrt{3{,}6903}=1{,}9210$ e $\sqrt{1{,}6427}=1{,}2817$), e o elemento fora da diagonal, a covariância. A correlação entre as disciplinas é
\[
r=\frac{c_{12}}{\sqrt{c_{11}\,c_{22}}}=\frac{1{,}4203}{1{,}9210\times 1{,}2817}=\frac{1{,}4203}{2{,}4622}\approx 0{,}577.
\]
A Figura~\ref{fig:corr} apresenta a matriz de correlação $R\in\R^{6\times 6}$ de $N$. Ela tem estrutura em blocos: dentro de cada disciplina as correlações entre períodos são altas (0,81 a 0,90 em Matemática; 0,82 a 0,89 em Português), enquanto entre disciplinas ficam entre 0,48 e 0,59. Em termos práticos, a nota do 2º período é um bom sinal antecipado da nota final ($r=0{,}90$ em Matemática), o que permite identificar estudantes em risco antes do fim do ano. Já a correlação moderada entre disciplinas indica que dificuldades em Matemática não são apenas reflexo de um desempenho geral baixo: há um componente específico da disciplina.

\begin{figure}[htbp]
\centering
\caption{Matriz de correlação $R=D^{-1/2}CD^{-1/2}$ das seis avaliações}
\label{fig:corr}
\includegraphics[width=0.46\textwidth]{graficos/fig3_matriz_correlacao.png}
\fonte{elaborada pelos autores (2026).}
\end{figure}

\subsection{Comparação entre grupos}

Para a variável escola, $E^TE=\operatorname{diag}(342;\ 40)$ e $G=(E^TE)^{-1}E^TF$ resultou em médias de 5,35 e 6,23 (Gabriel Pereira) contra 4,99 e 5,53 (Mousinho da Silveira), em Matemática e Português. A Figura~\ref{fig:grupos} mostra os demais grupos. O tempo de estudo tem relação positiva com o rendimento: em Matemática a média sobe de 5,13 (menos de 2 h semanais) para 5,84 (mais de 10 h). O histórico de reprovações separa os grupos com mais força: em Matemática, a média cai de 5,67 (nenhuma reprovação) para 4,02 (uma), 3,05 (duas) e 3,02 (três ou mais).

\begin{figure}[htbp]
\centering
\caption{Médias por grupo calculadas por $G=(E^TE)^{-1}E^TF$}
\label{fig:grupos}
\includegraphics[width=0.74\textwidth]{graficos/fig5_medias_por_grupo.png}
\fonte{elaborada pelos autores (2026).}
\end{figure}

\subsection{Análise crítica e robustez}

\textbf{Sensibilidade aos pesos.} A Tabela~\ref{tab:sens} mostra o efeito de três vetores $w$ calculados de uma só vez por $N_{\text{Mat}}W_c$. As médias variam menos de 0,05 ponto e as taxas de aprovação, no máximo 1,1 ponto percentual. Isso decorre da alta correlação entre períodos: quando as colunas de $N_{\text{Mat}}$ são parecidas, qualquer média ponderada delas é parecida. Logo, a conclusão de que Matemática é mais difícil não é um artefato da escolha dos pesos.

\begin{table}[htbp]
\centering
\caption{Sensibilidade dos indicadores ao vetor de pesos dos períodos}
\label{tab:sens}
\begin{singlespace}\small
\begin{tabular}{lcccc}
\toprule
 & \multicolumn{2}{c}{Matemática} & \multicolumn{2}{c}{Português}\\
\cmidrule(lr){2-3}\cmidrule(lr){4-5}
Vetor $w$ & Média & Aprovação & Média & Aprovação\\
\midrule
$(1/3;\ 1/3;\ 1/3)$ (uniforme) & 5,33 & 57,3\% & 6,14 & 85,9\%\\
$(0{,}3;\ 0{,}3;\ 0{,}4)$ (adotado) & 5,31 & 57,1\% & 6,16 & 85,9\%\\
$(0{,}2;\ 0{,}3;\ 0{,}5)$ (ênfase final) & 5,29 & 56,8\% & 6,18 & 84,8\%\\
\bottomrule
\end{tabular}
\end{singlespace}
\fonte{elaborada pelos autores a partir dos resultados do código (2026).}
\end{table}

\textbf{Notas zero.} Há 39 estudantes com zero no 3º período de Matemática e 5 em Português, quase sempre após notas positivas nos períodos anteriores, o que indica ausência na avaliação final, e não nota obtida. Excluindo os 42 estudantes nessa situação ($m=340$), a média de Matemática sobe de 5,31 para 5,72 e a aprovação de 57,1\% para 63,8\%; em Português, de 6,16 para 6,28 e de 85,9\% para 88,8\%. A diferença entre as disciplinas diminui, mas continua clara. O resultado mostra uma limitação do modelo linear: a matriz trata todo número como medida de desempenho, e cabe ao analista distinguir ``zero'' de ``dado ausente''.

\textbf{Limitações.} (i) A identificação dos estudantes comuns às duas bases é aproximada, pois depende de 13 atributos e não de um código individual. (ii) Os dados são de duas escolas de Portugal em 2005/2006 e não podem ser generalizados para outros contextos. (iii) As comparações entre grupos são descritivas: mostram associação, e não causa; grupos pequenos (por exemplo, $n=11$ estudantes com duas reprovações) produzem médias pouco estáveis. (iv) Os pesos $w$ e $p$ são parâmetros adotados, e não critérios oficiais.

% ================================ CONCLUSÃO ================================
\section{Conclusão}

O trabalho respondeu ao problema de pesquisa mostrando que os principais indicadores de desempenho acadêmico podem ser escritos e calculados como operações matriciais. Com a matriz de notas $N\in\R^{382\times 6}$, a multiplicação por escalar ajustou a escala; a soma de matrizes produziu médias e evoluções; o produto matriz-vetor, com $w$, e o produto por uma matriz bloco-diagonal calcularam as notas ponderadas de todos os estudantes de uma vez; a transposição com o vetor $\mathbf{1}$ gerou médias e taxas de aprovação; a matriz $X_c^TX_c$ revelou a variância e a correlação; e $(E^TE)^{-1}E^TF$ comparou grupos.

Os resultados indicam que Matemática é a disciplina de maior dificuldade (média 5,31, aprovação de 57,1\% e queda ao longo do ano), que o desempenho entre disciplinas é apenas moderadamente correlacionado, que a nota do 2º período antecipa bem a nota final e que o histórico de reprovações é o fator mais associado ao baixo rendimento. Esses achados sugerem ações concretas, como reforço de Matemática a partir do 2º período e acompanhamento de estudantes com reprovações anteriores. Como trabalhos futuros, propõe-se aplicar o mesmo modelo a dados de uma instituição brasileira, com mais disciplinas, e incluir técnicas como a decomposição em valores singulares para identificar perfis de estudantes.

% =============================== REFERÊNCIAS ===============================
\section*{REFERÊNCIAS}
\begin{singlespace}
\begin{flushleft}
\noindent ANTON, H.; RORRES, C. \textbf{Álgebra linear com aplicações}. 10. ed. Porto Alegre: Bookman, 2012.

\vspace{6pt}
\noindent BOLDRINI, J. L. \textit{et al}. \textbf{Álgebra linear}. 3. ed. São Paulo: Harbra, 1980.

\vspace{6pt}
\noindent CORTEZ, P.; SILVA, A. M. G. Using data mining to predict secondary school student performance. \textit{In}: BRITO, A.; TEIXEIRA, J. (ed.). \textbf{Proceedings of 5th Annual Future Business Technology Conference}. Porto: EUROSIS-ETI, 2008. p. 5-12.

\vspace{6pt}
\noindent CORTEZ, P. \textbf{Student performance}. Irvine: UCI Machine Learning Repository, 2014. Conjunto de dados. Disponível em: \url{https://archive.ics.uci.edu/dataset/320/student+performance}. Acesso em: 23 nov. 2026.

\vspace{6pt}
\noindent HARRIS, C. R. \textit{et al}. Array programming with NumPy. \textbf{Nature}, London, v. 585, p. 357-362, 2020. DOI: 10.1038/s41586-020-2649-2.

\vspace{6pt}
\noindent HUNTER, J. D. Matplotlib: a 2D graphics environment. \textbf{Computing in Science \& Engineering}, v. 9, n. 3, p. 90-95, 2007. DOI: 10.1109/MCSE.2007.55.

\vspace{6pt}
\noindent JOHNSON, R. A.; WICHERN, D. W. \textbf{Applied multivariate statistical analysis}. 6. ed. Upper Saddle River: Pearson Prentice Hall, 2007.

\vspace{6pt}
\noindent LAY, D. C. \textbf{Álgebra linear e suas aplicações}. 4. ed. Rio de Janeiro: LTC, 2013.

\vspace{6pt}
\noindent McKINNEY, W. Data structures for statistical computing in Python. \textit{In}: PYTHON IN SCIENCE CONFERENCE, 9., 2010, Austin. \textbf{Proceedings} [...]. Austin: SciPy, 2010. p. 56-61. DOI: 10.25080/Majora-92bf1922-00a.

\vspace{6pt}
\noindent ROMERO, C.; VENTURA, S. Educational data mining: a review of the state of the art. \textbf{IEEE Transactions on Systems, Man, and Cybernetics, Part C}, v. 40, n. 6, p. 601-618, 2010. DOI: 10.1109/TSMCC.2010.2053532.

\vspace{6pt}
\noindent STEINBRUCH, A.; WINTERLE, P. \textbf{Álgebra linear}. 2. ed. São Paulo: Pearson Makron Books, 1987.

\vspace{6pt}
\noindent STRANG, G. \textbf{Introduction to linear algebra}. 5. ed. Wellesley: Wellesley-Cambridge Press, 2016.
\end{flushleft}

\vspace{8pt}
\noindent\textbf{NOTA SOBRE O USO DE INTELIGÊNCIA ARTIFICIAL GENERATIVA}\par\vspace{3pt}
\noindent Os autores utilizaram o assistente de IA Claude (Anthropic) como apoio na estruturação do código, na redação preliminar do texto e na formatação. Todos os cálculos foram conferidos pela execução do código e por testes automatizados; a demonstração passo a passo foi refeita à mão; e todas as referências foram verificadas pelos autores. A responsabilidade pelo conteúdo é integralmente dos autores.
\end{singlespace}

\end{document}
